\section{Introduction}\label{sec:introduction}

Let $(X,B)$ be a projective log canonical pair over $\C$, and let
$f:X\to Z$ be a contraction: a surjective projective morphism of normal
projective varieties with $f_*\OO_X=\OO_Z$.  We consider the case
\[
 K_X+B\sim_\Q f^*D,
\]
where $D$ is a rational Cartier divisor on $Z$.  The adjoint of the
generic fibre is then rationally linearly trivial.  The canonical bundle
formula separates the contribution of singular fibres from a moduli
part on the base.  Its qualitative form places that moduli part on a
birational model as a nef rational Cartier divisor.  For effective
applications, one needs a Cartier multiple controlled by the dimension
and boundary coefficients, rather than by the individual fibration.

This is a question about a chosen rational presentation.  Replacing a
base divisor by an arbitrary rationally linearly equivalent divisor
changes the moduli part by a rational principal divisor and can introduce
arbitrarily large denominators.  We prove that, when $\dim X\leq4$ and the
coefficients of $B$ belong to a fixed finite rational set, one can choose
an exact pullback presentation whose moduli part has a uniform Cartier denominator.

We recall the base data in the statement.  For the given $f$, the
\emph{discriminant} $B_Z$ has coefficient $1-t_P$ at a prime divisor
$P\subset Z$, where $t_P$ is the log canonical threshold of $f^*P$ over
the generic point of $P$.  This is computed on a neighbourhood of that
point where $P$ is Cartier.  The moduli data form a rational b-divisor
$\mathbf M$.  On a birational model $q:W\to Z$ where it descends, it is
represented by a rational Cartier divisor $M_W$ and its pullbacks to
higher models; its trace on $Z$ is $M_Z=q_*M_W$.  The b-divisor is
\emph{b-nef} if $M_W$ is nef, and $p\mathbf M$ is \emph{b-Cartier} if
$pM_W$ is Cartier.  A generalized pair $(Z,B_Z+M_Z)$ has
$K_Z+B_Z+M_Z$ rational Cartier, and its crepant boundaries are defined by
\[
 K_W+B_W+M_W=q^*(K_Z+B_Z+M_Z)
\]
and by the same rule on higher models.  It is generalized lc if all
coefficients of these boundaries are at most one, and generalized klt
if they are strictly less than one.  The crepant boundaries on higher
models may have negative coefficients.

\begin{theorem}[Uniform relative denominators]\label{thm:relative-denominators}
Fix a finite set $\Phi\subset[0,1]\cap\Q$ and an integer $1\leq d\leq4$.
There are positive integers $p_0\mid p$, with $p_0$ clearing $\Phi$, and a
rational DCC set $\mathcal B\subset[0,1]$, depending only on $d$ and
$\Phi$, with the following property. Suppose $(X,B)$ is a projective lc
pair of dimension $d$ with coefficients in $\Phi$, and
\[
 f:X\longrightarrow Z,\qquad \dim Z>0,\qquad K_X+B\sim_\Q f^*D,
\]
where $f$ is a contraction and $D$ is a rational Cartier divisor. There
exist $\psi\in\C(X)^*$ and $D_Z\sim_\Q D$ such that
\begin{equation}\label{eq:relative-exact-presentation}
 K_X+B+\frac{1}{p_0}\Div(\psi)=f^*D_Z,
 \qquad D_Z=K_Z+B_Z+M_Z.
\end{equation}
Here $B_Z$ is effective with coefficients in $\mathcal B$, the pair
$(Z,B_Z+M_Z)$ is generalized lc, and $\mathbf M$ is b-nef with
$p\mathbf M$ b-Cartier. If $(X,B)$ is klt, the generalized pair is
generalized klt.
\end{theorem}

The first equality in~\eqref{eq:relative-exact-presentation} is an
equality of actual rational divisors, for compatible choices of canonical
divisors and the displayed rational function.  It selects the
representative $D_Z\sim_\Q D$ to which the moduli bound applies.  The two
integers have distinct roles: $p_0$ gives the principal correction and,
on the generic fibre, a trivialization in that degree; $p$ clears the
moduli trace on a model determining $\mathbf M$.  The discriminant is
controlled instead by the fixed DCC set $\mathcal B$.

\subsection{Denominators, positivity, and earlier work}

Fujino and Mori bound the denominator of the semistable divisor
accompanying $bK_X$ in their canonical bundle formula, where $b$ is the
least nonvanishing pluricanonical degree of the generic fibre.  Their
bound uses the middle Betti number of a smooth model of a cyclic cover of
the fibre of Kodaira dimension zero.  The proof reduces to a curve and
controls a finite character through this cohomology
\cite[Theorem~3.1 and Sections~3.3, 3.6--3.8]{FujinoMori2000}.  In their klt log Iitaka
setting, Todorov and Xu obtain a bound for the moduli denominator in terms
of the boundary coefficients in relative dimension two, with arbitrary
total dimension \cite[Theorem~1.3]{TodorovXu2009}.  Their proof first
controls surface indices and the second Betti number of a resolution of
the cyclic cover, then applies the Fujino--Mori character method.

When the generic fibre is a rational curve, Floris bounds a common integer
clearing all moduli coefficients in terms of the fibre Cartier index.  Her
examples rule out every polynomial bound for such an integer as the index
varies \cite[Theorem~1.6]{Floris2013}.  For the lc
total spaces of dimension at most four considered here, the curve
reduction produces reduced special fibres of dimension at most three.
The index of the whole reduced fibre supplies the source of uniformity
in our proof.

Qualitative moduli theory supplies a different part of the argument.
Ambro establishes rational b-Cartier descent and positivity in the
klt-trivial setting \cite{AmbroBoundary2004,AmbroModuli2005}, and Fujino
and Gongyo extend b-nefness to lc-trivial fibrations
\cite[Theorem~3.6]{FujinoGongyoModuli}.  These results put the moduli data
on a suitable base model and make the determining divisor nef; their
Cartier multiple may depend on the fibration.  We use this descent and
nefness after fixing the exact rational trivialization and verifying the
required rank-one condition, and obtain a common multiple by a separate
special-fibre calculation.  The preprint
of Bakker, Filipazzi, Mauri and Tsimerman proves qualitative
b-semiampleness for lc-trivial fibrations with generically effective
boundary \cite[Theorem~1.5]{BakkerFilipazziMauriTsimerman2025}.  Their
Section~7.2.4 discusses the general effective question and known special
cases, where one seeks a uniform multiple that is base point free on a
determining model.  A bound for the Cartier denominator makes a multiple
Cartier on such a model; base point freeness is an additional
condition, and is not used here.

The argument on a reducible fibre also has an established ancestry.
Fujino's admissible and preadmissible sections organize pluricanonical
descent across the conductor
\cite[Definition~4.1 and Lemma~4.2]{Fujino2000}.  Gongyo uses this
framework in the abundance theorem for numerically trivial slc adjoints
\cite[Theorem~1.5 and Section~5]{Gongyo2013}.  Jiang and Liu supply the
uniform index theorem for projective slc log Calabi--Yau pairs through
dimension three that we use below
\cite[Corollary~1.6, arXiv version~1]{JiangLiu2021}.  The relevant output
is one trivialization on the slc pair itself, with its conductor gluing.
It is this global form that permits the cyclic-character comparison in
the present proof.

\subsection{The route to the relative bound}

Section~\ref{sec:generic} applies the lower-dimensional index theorem to
the geometric generic fibre, whose dimension is at most three, and
descends its principalization by Hilbert~90 in the same degree $p_0$.
This fixes the exact presentation~\eqref{eq:relative-exact-presentation}.
Threshold ACC supplies the DCC set for $B_Z$.  After applying the
qualitative canonical bundle formula and passing to a smooth determining
model $W\to Z$, it remains to control one moduli coefficient at a time.
For a prime $P\subset W$, let $\alpha$ be the coefficient of the pullback
of $D_Z$, and let $t_P$ be the threshold computed from the crepant sub-pair
on a resolved diagram over $W$.  The coefficient of $M_W$ differs from
$\alpha+t_P$ by an integer.  Section~\ref{sec:curve} takes a transverse
curve slice and runs a relative dlt MMP to obtain an exact local identity
on a dlt model.
Adjunction in Section~\ref{sec:residue} then gives an slc pair $(T,B_T)$
on the whole connected reduced special fibre, with $\dim T\leq3$ and
$K_T+B_T\sim_\Q0$.

The decisive uniform input is a form on this entire fibre.  Global ACC
places the different coefficients in a fixed finite set.  The Jiang--Liu
index theorem then gives a nowhere-vanishing log pluricanonical form $u$ on $T$ in one even
multiple $p$ of $p_0$.  The local residue test compares its pullback with
the residues of an ambient form on a normalized cyclic cover; that
ambient form carries the character recording the remaining denominator.
The comparison may use a larger, pair-dependent degree for ambient
adjunction, but it yields one scalar on each normalized component of the
connected reduced covering fibre $T_Y$.  In the uniform degree $p$, the
local conductor calculation shows that each tuple has matching
nonzero next residues at a generic crossing of two components.  Their
scalars therefore agree across intersections and hence on all of $T_Y$.
The ambient tuple is therefore one scalar multiple of the invariant
pullback of $u$, even
when the cyclic group permutes components.  Its character is trivial in
degree $p$, giving $p(\alpha+t_P)\in\Z$ and the required Cartier bound.

\subsection{Effective systems and torsion}

When $D$ is big, Proposition~\ref{prop:effective-base} gives an integer
$m=m(d,\Phi)$ such that, for every positive multiple $l$ of $m$, the
complete rounded divisorial system
\[
 \left|\left\lfloor l(K_X+B)\right\rfloor\right|
\]
is nonempty and its section ratios generate exactly
$\C(Z)\subset\C(X)$.  Here the sheaf is the rank-one reflexive divisorial
sheaf.  Equality of the embedded fields recovers the finite part of the
fibration as well as the dimension of its image.

For every positive $l$ divisible by $p_0$, Section~\ref{sec:effective}
uses the exact presentation to identify the full section spaces inside
$\C(X)$:
\[
 H^0\!\left(X,\OO_X(\lfloor l(K_X+B)\rfloor)\right)
 =\psi^{l/p_0}f^*H^0\!\left(Z,\OO_Z(\lfloor lD_Z\rfloor)\right),
\]
where $f^*$ means pullback of rational functions.  The valuation argument
works for reflexive sheaves without requiring $l(K_X+B)$ to be Cartier.
Birkar--Zhang's polarized effective birationality theorem applies to the
resulting big lc adjoint on $W$, with fixed DCC boundary coefficients and
$pM_W$ nef Cartier \cite[Theorem~1.3]{BirkarZhang2016}.  Choosing $m$
divisible by $p_0$ and their uniform degrees, the equality above transfers
the full base field to every required source system: the common factor
cancels in ratios.  Their separate Iitaka-fibration bound
\cite[Theorem~1.2]{BirkarZhang2016} retains dependence on a fibre
nonvanishing degree and a Betti number of a cyclic fibre cover.

Proposition~\ref{prop:rc-torsion} is a separate torsion result.  Fix
$1\leq e\leq4$, a rational DCC set $I\subset[0,1]$, and a positive
integer $p$.  For projective generalized klt pairs $(Z,B_Z+M_Z)$ of
dimension $e$, with boundary coefficients in $I$, b-nef rational data
$\mathbf M$ for which $p\mathbf M$ is b-Cartier, and a rationally
connected smooth projective resolution of $Z$, it gives one integer
$\ell=\ell(e,I,p)$ such that
\[
 K_Z+B_Z+M_Z\sim_\Q0
 \quad\Longrightarrow\quad
 \ell(K_Z+B_Z+M_Z)\text{ is an integral principal divisor}.
\]
Using suitable extractions, global ACC places the boundary coefficients
in a fixed finite set and gives a uniform positive lower bound for
generalized log discrepancies.  Birkar's boundedness theorem
\cite[Theorem~1.7]{BirkarBoundedCY} bounds the varieties up to isomorphism
in codimension one.  The proof combines their smooth-locus topology with
Kummer theory and finite generation of the divisor class group $\Cl(Z)$
from the rationally connected resolution to give a uniform exponent for
torsion in $\Cl(Z)$.
A separate uniform multiple clears the
coefficients of the actual adjoint before that torsion bound is applied.
Birkar's ordinary-pair index consequence, observed by Totaro
\cite[Corollary~1.8 and Section~9.5]{BirkarBoundedCY}, is a close predecessor.

Corollary~\ref{cor:relative-rc-torsion} combines this result with
Theorem~\ref{thm:relative-denominators}.  A projective klt fourfold pair with fixed finite
rational boundary coefficients and $K_X+B\sim_\Q0$ has a uniform integral
principal multiple of $K_X+B$ whenever it admits a contraction to a
positive-dimensional base with a rationally connected smooth projective
resolution.  The statement includes $B=0$.
